\section{Poisson 计数、Exponential 等待与 Gamma 时间}
\label{sec:05-Probability/06-Common-Distributions/03}

你在监控一台服务器。可以问两个不同方向的问题：

\begin{itemize}
\tightlist
\item
  往后看一个小时，预计会收到多少个请求？——固定窗口，问计数。
\item
  从现在开始等，第 \(5\) 个请求什么时候到？——固定事件数，问时间。
\end{itemize}

同一个到达过程，两种观察角度。Poisson 分布回答计数问题；Exponential 分布回答等待第一个事件的时间问题；Gamma 分布回答等满 \(r\) 个事件的总时间问题。三者不是孤立的三个公式，是同一过程的三个投影。

\subsection{先把时间切碎：稀有事件作为大量小机会的叠加}

考虑一个长度为 \(t\) 的时间段。把它切成 \(n\) 个等分小格，每格极短，短到同一格里发生两次事件的概率可以忽略。每个小格里事件要么发生要么不发生。如果过程足够稳定——不相交区间的计数近似独立，小格事件概率大致恒定——那么整个 \(t\) 内的计数看起来就是一个 Binomial\((n,p)\)。

但这里的 \(n\) 和 \(p\) 有特殊关系。把格子越切越细（\(n\to\infty\)），每格的事件概率必然越来越小（\(p\to0\)），而平均计数 \(np\) 应该趋向一个有限值 \(\mu=\lambda t\)。这就是稀有事件 Binomial 的极限。

用一个具体数字感受这个过程。假设平均每小时 \(\lambda=3\) 个请求。你观测 \(t=1\) 小时。如果把一小时切成 \(n=60\) 个一分钟小格，每格概率 \(p=3/60=0.05\)，计数近似为 Binomial\((60,0.05)\)。切成 \(n=3600\) 个一秒小格，\(p=3/3600\approx0.000833\)，近似为 Binomial\((3600,0.000833)\)。两种切法的期望都是 \(3\)，但越细的切分越接近真实的 Poisson\((3)\)。到了极限，Binomail 的繁琐组合数被 Poisson 的简洁形式取代。

\subsection{从 Binomial 极限长出来的 Poisson}

固定 \(\mu=np\)。把 \(p=\mu/n\) 代入 Binomial PMF 并展开：

\[
\begin{aligned}
\binom{n}{k}p^k(1-p)^{n-k}
&=\frac{n(n-1)\cdots(n-k+1)}{k!}
\left(\frac{\mu}{n}\right)^k
\left(1-\frac{\mu}{n}\right)^{n-k}\\[2pt]
&=\frac{\mu^k}{k!}\cdot
\frac{n(n-1)\cdots(n-k+1)}{n^k}\cdot
\left(1-\frac{\mu}{n}\right)^{n}\cdot
\left(1-\frac{\mu}{n}\right)^{-k}.
\end{aligned}
\]

固定 \(k\)，令 \(n\to\infty\)。第二项趋近 \(1\)（分子分母都是 \(k\) 个因子的乘积，各因子比趋近 \(1\)）。第三项是 \((1-\mu/n)^n\)，标准极限 \(e^{-\mu}\)。第四项 \((1-\mu/n)^{-k}\to 1\)。乘在一起：

\[
\lim_{n\to\infty}\binom{n}{k}p^k(1-p)^{n-k}
=e^{-\mu}\frac{\mu^k}{k!}.
\]

这就是 Poisson 分布，记为 \(\mathrm{Poisson}(\mu)\)：

\[
P(N(t)=k)=e^{-\lambda t}\frac{(\lambda t)^k}{k!},\qquad k=0,1,2,\ldots
\]

参数 \(\lambda\) 是速率——单位时间内事件的平均发生次数。\(\mu=\lambda t\) 是该窗口的平均计数。

把这个极限当成一个数值对账来验证信用。取 \(\mu=3\)，\(n=100\)，\(p=0.03\)。Binomial\((100,0.03)\) 的 \(P(X=2)\) 是 \(\binom{100}{2}(0.03)^2(0.97)^{98}\approx0.225\)。Poisson\((3)\) 的 \(P(N=2)\) 是 \(e^{-3}\cdot 3^2/2!\approx0.224\)。两者只差千分之一量级——用极限代替精确 Binomial 的误差在 \(n=100\) 时就已经极小。这解释了 Poisson 为什么在计数数据建模中如此普及：大量机会加微小概率的组合在现实世界中频繁出现。

\subsection{\"均值等于方差\"不是巧合，是可以追查的诊断信号}

Poisson 的均值和方差相等，都等于 \(\mu=\lambda t\)。从定义直接验证：

\[
\begin{aligned}
\E[N]&=e^{-\mu}\sum_{k=1}^{\infty}\frac{\mu^k}{(k-1)!}=\mu,\\[2pt]
\E[N(N-1)]&=e^{-\mu}\sum_{k=2}^{\infty}\frac{\mu^k}{(k-2)!}=\mu^2.
\end{aligned}
\]

于是 \(\Var(N)=\E[N(N-1)]+\E[N]-(\E[N])^2=\mu^2+\mu-\mu^2=\mu\)。

这个等方差性质给了你一个很实用的诊断工具。拿到一组计数数据，如果样本方差显著大于样本均值（过离散），往往意味着速率不是常数——可能有时段密集、有时段稀疏，或者个体之间的速率本身有差异。比如你统计每个路口每周的事故数，有的路口在商圈、有的在郊区，把它们混在一起，方差会远超均值。如果样本方差显著小于均值（欠离散），可能存在容量限制或事件间的排斥效应——比如一个停车场每小时进入的车辆数，车位满了就进不去了。两种偏离都在告诉你：光用一个 \(\lambda\) 描述整个过程是不够的。

两个独立 Poisson 变量的和仍然是 Poisson：

\[
X\sim\mathrm{Poisson}(\lambda_1),\quad Y\sim\mathrm{Poisson}(\lambda_2)
\;\Longrightarrow\;
X+Y\sim\mathrm{Poisson}(\lambda_1+\lambda_2).
\]

这对应把两个独立的事件流叠在一起——两个独立的请求源，合并后的请求到达仍是 Poisson，速率相加。反方向的分解也存在：若总的计数是 Poisson\((\lambda)\)，而每个事件独立地以概率 \(p\) 被标记为``类型 A''，则类型 A 的计数服从 Poisson\((\lambda p)\)。这个性质叫 Poisson 的随机稀释，在分类计数问题中极其顺手。

\subsection{等待第一个事件：从\"没有事件\"反推等待时间}

换一个问法：第一个事件要等多久？记 \(T_1\) 为等待时间。``\(T_1>t\)''等价于``前 \(t\) 时间内一个事件都没有''。而 \(N(t)=0\) 的概率由 Poisson 的 \(k=0\) 项直接给出：

\[
P(T_1>t)=P(N(t)=0)=e^{-\lambda t}.
\]

有了生存函数，CDF 和密度就自然出来：

\[
F_{T_1}(t)=1-e^{-\lambda t},\qquad
f_{T_1}(t)=\lambda e^{-\lambda t},\qquad t\ge0.
\]

这就是 \(\mathrm{Exponential}(\lambda)\)——速率参数化。均值是 \(1/\lambda\)，方差是 \(1/\lambda^2\)。

用服务器的数字把直觉钉牢。\(\lambda=3\) 每小时，则平均等待时间是 \(1/3\) 小时 = \(20\) 分钟。但等待的波动很大——标准差也是 \(20\) 分钟。也就是说，有相当一部分时候你需要等 \(40\) 分钟以上，而有时可能几分钟就来。Exponential 的变异系数（标准差比均值）恒为 \(1\)：波动和中心同等量级，这一点在所有分布中是一个相当特别的性质。

有些教材用尺度参数 \(\theta=1/\lambda\)，密度写成 \((1/\theta)e^{-t/\theta}\)，均值就是 \(\theta\)。同一个分布、同一个符号、完全不同的参数角色——拿到公式先确认是 rate 还是 scale，否则算出的均值会互为倒数。

Exponential 最著名的性质：无记忆性。

\[
P(T>s+t\mid T>s)=\frac{e^{-\lambda(s+t)}}{e^{-\lambda s}}=e^{-\lambda t}=P(T>t).
\]

也就是说，无论你已经等了多久，剩余等待时间的分布和刚开始等时一模一样。这个性质在连续分布中只有 Exponential 满足（离散中则是 Geometric）。回到等公交的例子：你已经等了 \(15\) 分钟，剩下还要等的分布和一个刚来的人要等的分布完全相同——这在真实的公交场景中显然不对（公交车有大致时刻表），但对完全随机到达、没有调度的事件流，它是精确的。设备不老化、交通不分高峰低谷、排队的人不会午休。如果你的真实场景中这些因素存在，Exponential 就是一个需要被替换而非被修补的起点。

\subsection{等第 \(r\) 个事件：Gamma 是 Exponential 的叠加}

Poisson 过程的一个核心事实：相邻到达间隔是独立的 Exponential\((\lambda)\)。因此，等待第 \(r\) 个事件的总时间就是 \(r\) 个独立 Exponential\((\lambda)\) 之和：

\[
T_r = E_1+E_2+\cdots+E_r.
\]

这个总时间服从形状参数 \(r\)、速率参数 \(\lambda\) 的 Gamma 分布：

\[
f_{T_r}(t)=\frac{\lambda^r}{\Gamma(r)}\,t^{r-1}e^{-\lambda t},\qquad t>0.
\]

当 \(r\) 为正整数时，\(\Gamma(r)=(r-1)!\)。但 Gamma 函数的优雅在于它并不要求 \(r\) 是整数——任何正实数都可以作为形状参数，使 Gamma 族能灵活地拟合各种正值、右偏的等待或强度数据。非整数形状参数下，Gamma 不再能解释为整数个 Exponential 之和，但数学形式仍然合法。

均值和方差：

\[
\E[T_r]=\frac{r}{\lambda},\qquad
\Var(T_r)=\frac{r}{\lambda^2}.
\]

\(r=1\) 时退化为 Exponential。形状参数增大，密度从指数衰减的尖峰逐渐向钟形靠拢——你等的次数越多，总时间的相对波动越小。定量看：变异系数 \(\sqrt{\Var}/\E = 1/\sqrt{r}\)。等第 \(1\) 个时变异系数是 \(1\)，等第 \(10\) 个时降到约 \(0.32\)，等第 \(100\) 个时只有 \(0.1\)。这个收敛是中心极限定理在 Gamma 族上的具体映照：大量独立 Exponential 之和的分布趋向 Gaussian。

回到服务器。\(\lambda=3\)/小时。等第 \(5\) 个请求的平均时间 \(5/3\approx1.67\) 小时。标准差 \(\sqrt{5}/3\approx0.745\) 小时。单次等的波动巨大，但聚成批次后的可预测性明显提升。

注意 Gamma 的加法规则有条件：相互独立的 Gamma 变量只有速率参数相同时，加出来仍是 Gamma（形状参数相加）。速率不同则和不再留在 Gamma 族内。这个限制在使用 Gamma 混合模型时很重要。

\subsection{计数与等待是一句话的正反面}

这个等价关系虽然简单，但用好了可以绕过很多积分：

\[
\{T_r\le t\}\;\Longleftrightarrow\;\{N(t)\ge r\}.
\]

``第 \(r\) 个事件在时间 \(t\) 之前到达''和``在时间 \(t\) 内至少发生了 \(r\) 个事件''是同一件事的两种说法。因此

\[
P(T_r\le t)=P(N(t)\ge r)=1-\sum_{k=0}^{r-1}e^{-\lambda t}\frac{(\lambda t)^k}{k!}.
\]

右边是 Poisson 的尾部概率，计算时查 Poisson CDF 表比积分 Gamma 密度方便得多。这个关系也解释了为什么这三个分布总是成组出现：给定了 Poisson 过程的计数规律，等待时间的规律就被完全确定；反过来也一样。

\subsection{哪些情况会让你用错 Poisson}

Poisson 需要的假设很明确：事件在不相交区间内独立发生，速率恒定，瞬时不可能有两个事件同时发生。现实中偏离这三条的情况到处都是。

速率不是常数：早晚高峰的请求密度、季节性波动的故障率。这时需要非齐次 Poisson 过程——速率是时间的函数 \(\lambda(t)\)，计数不再是简单的 \(\lambda t\)，而变成速率函数在区间上的积分 \(\Lambda(t)=\int_0^t\lambda(s)\,ds\)，\(N(t)\sim\mathrm{Poisson}(\Lambda(t))\)。

事件之间不独立：一个故障引发连锁故障、一个顾客进店带动朋友进店。这种聚集效应使得方差远超均值，负二项分布或 Cox 过程（Poisson 的速率本身随机化）可能更合适。

瞬时多事件：批处理系统中一次可能到达多个任务。复合 Poisson 过程允许每次发生的事件数本身也是一个随机变量——总的计数可以拆成每次批量的叠加。

Exponential 无记忆性的背离更常见：设备越旧越容易坏（失效率递增→Weibull），或者刚修好的设备需要磨合期（失效率递减）。Gamma 的灵活性部分弥补了这一点，但真正需要的是检查数据是否支持无记忆假设，而不是默认如此。一个简单的诊断是画间隔时间的 Q-Q 图：若点偏离对角线呈弯曲趋势，无记忆性就值得怀疑。

Poisson--Exponential--Gamma 体系的美妙之处在于它们是一体三面。理解了计数和等待的等价关系，三个分布就不再是记忆列表中的三个独立条目，而是你手里的一套切换视角的工具。你可以从计数翻到等待，从等待翻到累积，每个角度回答一类具体问题——切换不需要重新建模，因为它们的根基是同一个过程。
